\chapter{Grafy Cayleya, działania na drzewach i teoria Bassa--Serre'a}
\label{ch:tg-bass-serre}

Dotychczas traktowaliśmy graf jako skończoną przestrzeń kombinatoryczną lub
jednowymiarowy kompleks CW. Teraz dopuszczamy także grafy nieskończone. Niektóre z nich,
jak graf Cayleya grupy względem skończonego zbioru generatorów,
są lokalnie skończone: każdy wierzchołek ma skończony stopień.
Natomiast drzewa pojawiające się w teorii Bassa--Serre'a mogą
mieć wierzchołki nieskończonego stopnia. Grafy nieskończone
występują naturalnie jako nakrycia kompleksów oraz przestrzenie,
na których działają grupy. Działania na drzewach pozwalają
odczytywać rozkłady grup fundamentalnych, a następnie rozkłady rozmaitości.

\section{Grafy Cayleya i prezentacje grup}
\label{sec:tg-cayley}

\begin{definicja}[Graf Cayleya]
Niech $G$ będzie grupą, a $S\subseteq G\setminus\{1\}$
skończonym, symetrycznym zbiorem, czyli $S=S^{-1}$.
\emph{Grafem Cayleya} $\operatorname{Cay}(G,S)$ nazywamy graf,
którego zbiorem wierzchołków jest $G$, a krawędzie łączą
$g$ z $gs$ dla wszystkich $g\in G$ i $s\in S$.
Każdą nieorientowaną krawędź $\{g,gs\}$ liczymy jednokrotnie,
także gdy $s=s^{-1}$. Ponieważ $1\notin S$,
graf nie ma pętli, a skracanie w grupie pokazuje,
że różne elementy $S$ dają różnych sąsiadów $g$.
Stopień każdego wierzchołka wynosi więc $|S|$. Gdy potrzebujemy zachować pełną informację o generatorach,
rozpatrujemy graf skierowany z łukami $g\to gs$
etykietowanymi przez $s$. W niezorientowanym grafie
utożsamiamy łuk $(g,s)$ z łukiem $(gs,s^{-1})$.
\end{definicja}

\begin{stwierdzenie}
Niech $G$ będzie grupą, a $S\subseteq G\setminus\{1\}$
skończonym zbiorem symetrycznym. Graf $\operatorname{Cay}(G,S)$
jest spójny wtedy i tylko wtedy, gdy $S$ generuje $G$.
\end{stwierdzenie}
\begin{proof}
Jeżeli $S$ generuje $G$, to dla dowolnego $g\in G$ istnieje
rozkład $g=s_1\cdots s_k$, gdzie $s_i\in S$. Kolejne
wierzchołki $1,s_1,s_1s_2,\ldots,g$ tworzą spacer od
jedności do $g$. Ponieważ z każdego $g$ otrzymujemy
spacer do $h$ przez translację spaceru od $1$ do
$g^{-1}h$, graf jest spójny. Odwrotnie, jeżeli istnieje
spacer $1=g_0,g_1,\ldots,g_k=g$, to dla każdego $i$
z definicji krawędzi $g_{i-1}^{-1}g_i\in S$.
Po wymnożeniu równości dostajemy
$g=(g_0^{-1}g_1)\cdots(g_{k-1}^{-1}g_k)$, a więc
każdy element należy do podgrupy generowanej przez $S$.
\end{proof}

\begin{przyklad}[Trzy rodziny grafów Cayleya]
Dla $G=\mathbb Z$ i $S=\{-1,1\}$ otrzymujemy
nieskończoną prostą, z wierzchołkami oznaczonymi liczbami całkowitymi.
Dla $G=\mathbb Z/n\mathbb Z$, $n\ge3$, oraz $S=\{-1,1\}$
otrzymujemy cykl $C_n$. Dla $G=\mathbb Z^2$ i
$S=\{(\pm1,0),(0,\pm1)\}$ powstaje kwadratowa sieć
płaszczyzny. W dwóch ostatnich przykładach odczytujemy
z etykietowanych zamkniętych spacerów odpowiednio relację $a^n=1$ i przemienność
generatorów odpowiadających przesunięciom poziomym
oraz pionowym.
\end{przyklad}

\begin{center}
\begin{tikzpicture}[line cap=round,>=Stealth,scale=.94]
 \foreach \x in {-2,-1,0,1,2} {
   \node[tgBlueVertex] (z\x) at (\x*.58,0) {};
   \node[font=\scriptsize,text=tgNavy] at (\x*.58,-.38) {$\x$};
 }
 \draw[tgEdge] (z-2)--(z-1)--(z0)--(z1)--(z2);
 \draw[tgEdge,dashed] (-1.70,0)--(-1.38,0) (1.38,0)--(1.70,0);
 \node[font=\small,text=tgNavy] at (0,.62) {$\operatorname{Cay}(\mathbb Z,\{\pm1\})$};
 \begin{scope}[shift={(4.45,0)}]
  \foreach \a in {0,60,...,300} {
    \node[tgTealVertex] (c\a) at (\a:1.02) {};
    \node[font=\scriptsize,text=tgNavy] at (\a:1.37) {$\pgfmathparse{int(\a/60)}\pgfmathresult$};
  }
  \foreach \a [evaluate=\a as \b using {int(mod(\a+60,360))}] in {0,60,...,300} {
    \draw[tgTree] (c\a)--(c\b);
  }
  \node[font=\small,text=tgNavy] at (0,-1.72) {$\operatorname{Cay}(\mathbb Z/6\mathbb Z,\{\pm1\})$};
 \end{scope}
\end{tikzpicture}
\end{center}

\begin{twierdzenie}[Graf Cayleya grupy wolnej]
Niech $F_r$, $r\ge1$, będzie grupą wolną
z ustaloną wolną bazą $a_1,\ldots,a_r$, a $S=\{a_1^{\pm1},\ldots,a_r^{\pm1}\}$.
Wtedy $\operatorname{Cay}(F_r,S)$ jest drzewem,
a każdy jego wierzchołek ma stopień $2r$
(również dla $r=1$, gdy otrzymujemy prostą).
\end{twierdzenie}
\begin{proof}
Z poprzedniego stwierdzenia graf jest spójny. Gdyby istniał
cykl prosty $g_0,g_1,\ldots,g_k=g_0$, $k\ge3$, to
kolejne etykiety krawędzi dawałyby relację
$s_1\cdots s_k=1$ w $F_r$. Ponieważ cykl nie zawraca
natychmiast tą samą krawędzią, słowo nie zawiera sąsiedniej
pary $a_i a_i^{-1}$ ani $a_i^{-1}a_i$; jest więc słowem
zredukowanym i niepustym. W grupie wolnej słowo takie
nie reprezentuje jedności (por. definicję grupy wolnej
w rozdziale \ref{ch:tg-topology}), co daje sprzeczność.
Spójny graf bez cykli jest drzewem. Dla każdego
$r\ge1$ różne generatory i ich odwrotności
dają $2r$ różnych sąsiadów wierzchołka,
gdyż w wolnej grupie żaden z tych elementów
nie jest równy innemu.
\end{proof}

\section{Działania grup na drzewach i ilorazy}
\label{sec:tg-tree-actions}

\begin{definicja}[Działanie bez inwersji]
Niech grupa $G$ działa przez automorfizmy na drzewie $T$.
Mówimy, że działanie jest \emph{bez inwersji}, jeżeli
żaden element $g\in G$ nie zamienia końców krawędzi,
pozostawiając tę krawędź jako zbiór niezmienniczy.
Stabilizatorem wierzchołka $v$ jest
$G_v=\{g\in G:gv=v\}$, a krawędzi $e$ stabilizator
$G_e=\{g\in G:ge=e\}$. Przy działaniu bez inwersji
dla krawędzi o końcach $v,w$ zachodzi
$G_e=G_v\cap G_w$: stabilizacja krawędzi wymusza
ustalenie obu końców, a ustalenie końców ustala
ich jedyną krawędź w drzewie.
\end{definicja}

\begin{stwierdzenie}[Wolne działania na drzewach]
Niech grupa $G$ działa swobodnie na wierzchołkach
niepustego drzewa $T$ i bez inwersji, przez automorfizmy grafu.
Realizacje grafów wyposażamy w topologię kompleksów CW,
także przy nieskończonych stopniach wierzchołków. Wtedy odwzorowanie
$T\to G\backslash T$ jest nakryciem grafów, $T$ jest
nakryciem uniwersalnym ilorazu, a grupa $G$ jest wolna.
Jeżeli iloraz jest skończonym grafem spójnym z $v$
wierzchołkami i $e$ krawędziami, to $G\cong F_{e-v+1}$.
\end{stwierdzenie}
\begin{proof}
Niech $q:T\to G\backslash T$ będzie odwzorowaniem
ilorazowym. Wierzchołki ilorazu to orbity wierzchołków,
a zorientowane krawędzie to orbity zorientowanych
krawędzi $T$. Odwrócenie orientacji indukuje
inwolucję bez punktów stałych na tych orbitach:
gdyby $e$ i $\bar e$ leżały w jednej orbicie,
pewien element grupy odwracałby krawędź.
Początek i koniec są dobrze określone jako orbity
odpowiednich końców. Ustalmy wierzchołek $v\in T$.
Dla dowolnej zorientowanej półkrawędzi ilorazu
rozpoczynającej się w $q(v)$ wybierzmy jej
reprezentanta $e$ o początku $gv$ dla pewnego
$g\in G$. Wtedy półkrawędź $g^{-1}e$
rozpoczyna się w $v$ i ma żądany obraz.
To dowodzi surjektywności odwzorowania gwiazd.

Gdy dwie półkrawędzie $e_1,e_2$ o początku $v$
mają ten sam obraz, istnieje $g\in G$,
który przenosi jedną na drugą. Ponieważ działanie
jest bez inwersji, $gv=v$, a zatem swoboda
wymusza $g=1$ i $e_1=e_2$.
Gwiazdy odwzorowują się więc bijektywnie.
W realizacji CW wybieramy po jednej orientacji
każdej krawędzi ilorazu i orientujemy krawędzie $T$
zgodnie z jej podniesieniami. Na odcinkach
parametryzujących krawędzie działanie grupy
przedłużamy liniowo; zachowuje ono parametr
w wybranej orientacji. Odwzorowanie $q$ jest
liniowe na każdej krawędzi i ciągłe w topologii CW.
Otwarte gwiazdy z ramionami krótszymi niż połowa
krawędzi są równomiernie nakryte: ich przeciwobrazy
są rozłącznymi gwiazdami, a bijekcja ramion daje
homeomorfizm każdej z nich na gwiazdę bazową.
Nawet przy nieskończonej liczbie ramion ciągłość
w obie strony wynika z topologii CW, w której
otwartość sprawdza się na każdej domkniętej krawędzi.
Wnętrza krawędzi są również równomiernie nakryte.

Drzewo jest spójne i po prostu spójne: każda pętla
ma zwarty obraz, a zwarty podzbiór kompleksu CW
jest zawarty w skończonym podkompleksie
(standardowa własność kompleksów CW).
Skończenie wiele krawędzi drzewa można zawrzeć
w skończonym poddrzewie, dołączając jedyne drogi
do ustalonego korzenia. Takie poddrzewo jest
kontraktowalne przez kolejne usuwanie liści
(por. dowód twierdzenia \ref{thm:tg-wedge}).
Zatem każda pętla jest ściągalna, a $q$ jest
nakryciem uniwersalnym.

Grupa przekształceń pokładowych tego nakrycia
jest dokładnie $G$: zawiera ona dane działanie,
a każde przekształcenie pokładowe jest jednoznacznie
wyznaczone obrazem jednego wierzchołka: dwa takie
przekształcenia zgodne w punkcie bazowym są
podniesieniami tego samego odwzorowania $q$,
więc z jednoznaczności podniesienia dróg i spójności
drogowej drzewa są zgodne wszędzie.
Ponieważ wierzchołki w tym samym włóknie
stanowią jedną orbitę $G$, taki obraz można
zrealizować odpowiednim elementem $G$.
Dostajemy $G\cong\pi_1(G\backslash T)$,
z dokładnością do zwykłego wyboru punktu bazowego
i konwencji lewego lub prawego działania.
Grupa fundamentalna spójnego grafu jest wolna
również w przypadku nieskończonym. Istotnie,
wybieramy drzewo rozpinające (istnieje na mocy
lematu Zorna, stosowanego do lasów zawierających
wszystkie wierzchołki grafu, uporządkowanych
przez inkluzję krawędzi: suma łańcucha lasów
nie zawiera cyklu, gdyż cykl ma skończenie wiele
krawędzi. Maksymalny las musi być spójny, ponieważ
w przeciwnym razie można dodać krawędź łączącą
dwie różne składowe, korzystając ze spójności grafu). Wybieramy jedną orientację każdej
krawędzi spoza drzewa. Generatory to pętle od
korzenia po drzewie do początku takiej krawędzi,
przez nią i z powrotem po drzewie.
Każda pętla i każda homotopia pętli ma obraz
zawarty w skończonym podkompleksie CW.
Dołączając do niego skończenie wiele dróg drzewowych
od korzenia do jego wierzchołków, otrzymujemy
skończony spójny graf, którego część drzewowa jest
drzewem rozpinającym. Twierdzenie
\ref{thm:tg-free-pi1} dowodzi zarówno generowania,
jak i braku relacji między wybranymi generatorami:
każda ewentualna relacja miałaby homotopię
w takim skończonym grafie. Dla skończonego ilorazu
liczba generatorów wynosi $e-v+1$.
\end{proof}

\section{Grafy grup i twierdzenie van Kampena}
\label{sec:tg-graphs-of-groups}

\begin{definicja}[Graf grup]
Niech $\Gamma$ będzie spójnym grafem, którego każdą
krawędź rozpatrujemy w dwóch przeciwnych orientacjach
$e$ oraz $\bar e$. \emph{Graf grup} składa się z grup
$G_v$ przypisanych wierzchołkom, grup $G_e=G_{\bar e}$
przypisanych krawędziom oraz monomorfizmów
$\iota_{e,-}:G_e\hookrightarrow G_{o(e)}$ i
$\iota_{e,+}:G_e\hookrightarrow G_{t(e)}$.
Zmiana orientacji krawędzi wymienia oba monomorfizmy:
$\iota_{\bar e,-}=\iota_{e,+}$ oraz
$\iota_{\bar e,+}=\iota_{e,-}$.
\end{definicja}

\begin{definicja}[Grupa fundamentalna grafu grup]
Ustalmy drzewo rozpinające $D$ grafu $\Gamma$.
Grupę fundamentalną grafu grup względem $D$
definiujemy przez prezentację
\[
 \left\langle *_{v\in V(\Gamma)}G_v,\ t_e
 \ \middle|\
 \begin{array}{l}
 t_{\bar e}=t_e^{-1},\\
 t_e=1\quad(e\in D),\\
 t_e\iota_{e,-}(c)t_e^{-1}=\iota_{e,+}(c)
       \quad(c\in G_e)
 \end{array}\right\rangle .
\]
Oznacza to iloraz iloczynu wolnego grup
wierzchołkowych i grupy wolnej na literach
$t_e$ przez podgrupę normalną generowaną
przez wypisane relacje. W szczególności zachowujemy
wszystkie relacje wewnętrzne grup wierzchołkowych.
Wybór jednej orientacji krawędzi pozwala usunąć
litery $t_{\bar e}$; relacje dla krawędzi drzewa
utożsamiają oba obrazy grupy krawędziowej.
Standardowe twierdzenie teorii Bassa--Serre'a
(zob. literatura na końcu rozdziału) zapewnia
niezależność typu izomorfizmu od $D$ i injektywność
naturalnych homomorfizmów grup wierzchołkowych
do grupy fundamentalnej. Tych dwóch faktów
nie wyprowadzamy z samej prezentacji.
\end{definicja}

\begin{twierdzenie}[Dwa podstawowe przypadki van Kampena]
\label{thm:tg-vk-amalgam-hnn}
\begin{enumerate}[label=(\roman*)]
\item Niech $X=A\cup B$ będzie łukowo spójną
przestrzenią, gdzie $A,B,A\cap B$ są otwarte,
łukowo spójne i zawierają wspólny punkt bazowy.
Załóżmy, że homomorfizmy
$\pi_1(A\cap B)\to\pi_1(A)$ i
$\pi_1(A\cap B)\to\pi_1(B)$ są injekcjami.
Wtedy
$\pi_1(X)\cong\pi_1(A)*_{\pi_1(A\cap B)}\pi_1(B)$.
\item Niech $Y,Z$ będą łukowo spójnymi przestrzeniami,
a $\iota_-,\iota_+\colon Z\to Y$ ciągłymi odwzorowaniami
indukującymi monomorfizmy na grupach fundamentalnych
(po wyborze zgodnych punktów bazowych). Jeśli przestrzeń
jest ilorazem powstałym przez dołączenie cylindra $Z\times[0,1]$ do $Y$
za pomocą $\iota_-$ i $\iota_+$ na jego końcach,
to jej grupa fundamentalna jest
rozszerzeniem HNN grupy $\pi_1(Y)$ po odpowiednio
utożsamionych obrazach $\pi_1(Z)$.
\end{enumerate}
\end{twierdzenie}
\begin{proof}
W przypadku (i) klasyczne twierdzenie van Kampena stwierdza, że
$\pi_1(X)$ jest ilorazem iloczynu wolnego
$\pi_1(A)*\pi_1(B)$ przez relacje utożsamiające
oba obrazy każdego elementu $\pi_1(A\cap B)$.
Ponieważ oba odwzorowania są injektywne, jest to
właśnie iloczyn wolny z amalgamacją. W przypadku (ii) przestrzenie $Y$ i $Z$ są łukowo spójne,
a iloraz powstaje ze $Y\sqcup(Z\times[0,1])$ przez
utożsamienie $(z,0)$ z $\iota_-(z)$ i $(z,1)$ z
$\iota_+(z)$; oba homomorfizmy $\pi_1(Z)\to\pi_1(Y)$
są injekcjami. Ustalmy $z_0\in Z$, $y_0\in Y$
oraz ścieżki w $Y$ od $y_0$ do obrazów $z_0$.
Oznaczmy te ścieżki przez $p_-,p_+$ i niech
$\gamma$ będzie drogą $s\mapsto(z_0,s)$ od końca
minus do końca plus. W tym dowodzie drogi
składamy od lewej do prawej. Homomorfizmy bazowe
oznaczamy przez $\iota_{\pm*}$, zatem
$\iota_{\pm*}([c])=[p_\pm\,\iota_\pm(c)\,p_\pm^{-1}]$.
Wybieramy literę
$t=[p_+\gamma^{-1}p_-^{-1}]$.

W przestrzeni ilorazowej bierzemy otwarty zbiór
$U$ złożony z $Y$ oraz kołnierzy cylindra
$Z\times([0,1/3)\cup(2/3,1])$ oraz zbiór
$V=Z\times(1/4,3/4)$. Otwartość wynika
z otwartości ich przeciwobrazów w rozłącznej sumie.
Zbiór $U$ deformacyjnie ściąga się na $Y$:
na $Y$ homotopia jest identycznością, a na dolnym
kołnierzu zastępujemy parametr $s$ przez $(1-u)s$,
na górnym przez $1-(1-u)(1-s)$, gdzie $u\in[0,1]$.
Definicje są zgodne z utożsamieniami na końcach
cylindra. Ciągłość homotopii wynika z topologii
ilorazowej: ograniczenie mapy ilorazowej do
przeciwobrazu otwartego $U$ jest ilorazowe, a jej
iloczyn z identycznością przedziału $[0,1]$
pozostaje ilorazowy (por. dowód lematu
\ref{lem:tg-collapse-CW}).
Zbiór $V$ deformacyjnie ściąga się na kopię $Z$
w połowie cylindra przez liniową zmianę parametru.
Przecięcie ma dokładnie dwie łukowo spójne składowe,
$Z\times(1/4,1/3)$ i $Z\times(2/3,3/4)$.
Grupoidowa wersja twierdzenia van Kampena
(z dwoma punktami bazowymi, po jednym w każdej
składowej przecięcia) daje prezentację z grupami
$\pi_1(Y)$ i $\pi_1(Z)$ oraz jedną dodatkową
literą odpowiadającą połączeniu obu punktów
przez $V$ i powrotowi przez $U$.
Relacje z pierwszej składowej eliminują grupę
$\pi_1(Z)$ przez utożsamienie jej z obrazem
$\iota_{-*}$, a relacje z drugiej składowej
utożsamiają jej sprzężony obraz z $\iota_{+*}$.
Kierunek sprzężenia sprawdzamy w cylindrze:
dla pętli $c:[0,1]\to Z$ o końcach $z_0$
kwadrat $(r,s)\mapsto[(c(r),s)]$ w ilorazie
ma dolny i górny bok odpowiadające
$\iota_-(c)$ i $\iota_+(c)$, a oba pionowe
boki odpowiadają $\gamma$.
Stąd droga $\gamma^{-1}\iota_-(c)\gamma$
jest homotopijna relatywnie do końców
z $\iota_+(c)$.
Stąd $t\iota_{-*}([c])t^{-1}=\iota_{+*}([c])$.
Dostajemy prezentację
\[
\langle \pi_1(Y),t\mid t\iota_{-*}(c)t^{-1}=\iota_{+*}(c)
\ \text{dla wszystkich }c\in\pi_1(Z)\rangle.
\]
Ta prezentacja definiuje rozszerzenie HNN.
Jeżeli zmienimy ścieżki bazowe na $p'_\pm$,
oznaczmy $u_\pm=[p'_\pm p_\pm^{-1}]$.
Nowe homomorfizmy są postaci
$\iota'_{\pm*}(c)=u_\pm\iota_{\pm*}(c)u_\pm^{-1}$,
a nowa litera to $t'=u_+tu_-^{-1}$.
Podstawienie tych wzorów przenosi relacje
jednej prezentacji na drugą; podstawienie
odwrotne daje izomorfizm. W szczególności nie można
pominąć nowego generatora $t$: homomorfizm na
$\mathbb Z$ przypisujący $t$ liczbę $1$, a całej
grupie $\pi_1(Y)$ liczbę $0$ jest dobrze określony,
ponieważ obie strony każdej relacji HNN mają
obraz zero.
\end{proof}

\begin{przyklad}[Iloczyn wolny i torus jako HNN]
Gdy graf grup jest odcinkiem, grupy na jego końcach to
$\mathbb Z=\langle a\rangle$ i
$\mathbb Z=\langle b\rangle$, a grupa krawędzi jest
trywialna, otrzymujemy $\mathbb Z*\mathbb Z\cong F_2$.
Gdy graf jest jedną pętlą, grupa wierzchołka i krawędzi
są równe $\mathbb Z$, a oba monomorfizmy grupy krawędzi
w grupę wierzchołka są identycznościami, otrzymujemy
\[
\langle a,t\mid tat^{-1}=a\rangle\cong\mathbb Z^2.
\]
Pierwszy przykład odpowiada sklejaniu dwóch okręgów
w punkcie, drugi jest przykładem cylindrycznego
sklejania prowadzącego do torusa.
\end{przyklad}

\begin{center}
\begin{tikzpicture}[>=Stealth,scale=1]
 \node[tgBlueVertex] (a) at (-2.5,0) {};
 \node[tgRedVertex] (b) at (-.5,0) {};
 \draw[tgEdge] (a)--node[above,font=\small] {$G_e$}(b);
 \node[below=5pt of a,font=\scriptsize] {$G_{v_1}$};
 \node[below=5pt of b,font=\scriptsize] {$G_{v_2}$};
 \node[tgTealVertex] (c) at (2.5,0) {};
 \draw[tgAccent] (c) .. controls (1.65,.95) and (3.35,.95) .. (c);
 \node[font=\scriptsize,text=tgNavy] at (2.5,1.16) {jedna krawędź pętlowa};
 \node[below=5pt of c,font=\scriptsize] {$G_v$};
 \node[text=tgNavy,font=\scriptsize] at (-1.5,-.80) {amalgamacja};
 \node[text=tgNavy,font=\scriptsize] at (2.5,-.80) {rozszerzenie HNN};
\end{tikzpicture}
\end{center}

\begin{uwaga}[Wymagane dane grafu grup]
Nie wystarcza znajomość abstrakcyjnego grafu ilorazowego
$G\backslash T$. Konstrukcja potrzebuje stabilizatorów
wierzchołków i krawędzi, homomorfizmów włączających oraz
wybranych elementów realizujących zmianę przedstawicieli orbit.
Dwa działania o jednakowym grafie ilorazowym mogą dawać
nieizomorficzne grupy fundamentalne.
Ponadto drzewo Bassa--Serre'a nie musi być lokalnie skończone:
stopnie wierzchołków zależą od indeksów podgrup krawędziowych.
\end{uwaga}

\begin{twierdzenie}[Fundamentalne twierdzenie Bassa--Serre'a]
Niech $G$ działa bez inwersji na drzewie $T$.
Po wybraniu przedstawicieli orbit wierzchołków i krawędzi,
dróg do tych przedstawicieli i odpowiednich elementów grupy,
działanie określa graf grup na grafie ilorazowym
$G\backslash T$. Jego grupami wierzchołków i krawędzi
są stabilizatory, a monomorfizmy krawędziowe pochodzą
z inkluzji stabilizatorów po potrzebnych sprzężeniach. Grupa fundamentalna tego
grafu grup jest izomorficzna z $G$. Odwrotnie, każdy graf
grup posiada drzewo Bassa--Serre'a, na którym jego grupa
fundamentalna działa bez inwersji, a otrzymany graf grup
jest izomorficzny z wyjściowym (po odpowiednich wyborach).
\end{twierdzenie}

\begin{uwaga}[Zakres dowodu]
Powyższy wynik jest bardziej ogólny niż twierdzenie van
Kampena o amalgamacji i HNN. Do pełnego dowodu trzeba
skonstruować normalną postać słów i drzewo Bassa--Serre'a
jako graf klas bocznych oraz sprawdzić, że nie ma w nim
cykli. Są to dodatkowe wyniki, które tutaj przyjmujemy
z monografii J.-P. Serre'a \emph{Trees}
(por. literatura na końcu rozdziału).
Dowiedzione wcześniej przypadki pokazują natomiast,
dlaczego w rozkładach topologicznych pojawiają się właśnie
iloczyny wolne z amalgamacją i rozszerzenia HNN.
\end{uwaga}

\begin{stwierdzenie}[Stopnie w drzewie Bassa--Serre'a]
Niech $\mathcal G$ będzie spójnym grafem grup o grafie bazowym
$\Gamma$, a $T$ jego drzewem Bassa--Serre'a. Dla
wierzchołka $\widetilde v\in T$ nad $v\in V(\Gamma)$
liczba półkrawędzi wychodzących z $\widetilde v$ wynosi
\[
\deg_T(\widetilde v)=
\sum_{e\colon o(e)=v}[G_v:\iota_{e,-}(G_e)],
\]
przy czym sumujemy po wszystkich zorientowanych
półkrawędziach grafu $\Gamma$ rozpoczynających się w $v$;
w pętli obie orientacje są osobnymi składnikami.
Sumy i indeksy mogą być nieskończone.
\end{stwierdzenie}
\begin{proof}
W konstrukcji drzewa Bassa--Serre'a wierzchołki nad $v$
są lewymi klasami bocznymi grupy $G_v$ w grupie
fundamentalnej grafu grup. Wystarczy więc rozpatrzyć
wierzchołek odpowiadający klasie $G_v$, którego
stabilizator to $G_v$ (pozostałe mają sprzężone
stabilizatory i tę samą liczbę incydencji).
W konstrukcji klas bocznych półkrawędź typu $e$
jest określona przez klasę $g\iota_{e,-}(G_e)$.
Przy przyjętej wcześniej konwencji
$t_e\iota_{e,-}(c)t_e^{-1}=\iota_{e,+}(c)$
jej początek to $gG_{o(e)}$, a koniec
$gt_e^{-1}G_{t(e)}$. Typy wierzchołków i krawędzi
zachowujemy jako część danych, nawet gdy
odpowiednie podgrupy są równe.
Jeżeli zmienimy $g$ na $g\iota_{e,-}(c)$,
początek pozostaje ten sam, a koniec też,
ponieważ $\iota_{e,-}(c)t_e^{-1}
=t_e^{-1}\iota_{e,+}(c)$. Odwrócona półkrawędź
ma typ $\bar e$ i klasę
$gt_e^{-1}\iota_{e,+}(G_e)$.
Ponowne odwrócenie daje pierwotną klasę,
gdyż $t_{\bar e}^{-1}=t_e$. Lewa translacja
przez grupę fundamentalną zachowuje te dane.
Fakt, że tak skonstruowany graf jest drzewem,
jest częścią przyjętego twierdzenia Bassa--Serre'a.
Dla początku $G_v$
mamy $g\in G_v$. Półkrawędzie tego typu odpowiadają
zatem bijektywnie lewym klasom bocznym
$G_v/\iota_{e,-}(G_e)$, rozumianym jako zbiór,
nie jako grupa ilorazowa. Ich liczba to podany indeks.
Zsumowanie po różnych orbitach krawędzi daje wzór.
\end{proof}
\begin{przyklad}
Dla iloczynu wolnego $C_2*C_3$ graf grup ma dwa wierzchołki
z grupami $C_2$ i $C_3$ oraz jedną krawędź z grupą trywialną.
Drzewo Bassa--Serre'a jest dwudzielne: stopnie wierzchołków
obu typów wynoszą odpowiednio $2$ i $3$. Natomiast
$\mathbb Z*\mathbb Z$ otrzymane z dwóch grup $\mathbb Z$
i trywialnej grupy krawędzi daje drzewo o nieskończonych
stopniach wierzchołków. Nie można więc zakładać lokalnej
skończoności drzew Bassa--Serre'a bez dodatkowych hipotez.
\end{przyklad}

\begin{przyklad}[Przypadek pętlowy: grupa Baumslaga--Solitara]
Rozważmy graf grup o jednym wierzchołku, jednej
krawędzi pętlowej i grupach wierzchołka oraz
krawędzi równych $\mathbb Z=\langle a\rangle$.
Niech $\iota_{e,-}$ będzie identycznością,
a $\iota_{e,+}$ mnożeniem przez dwa.
Wtedy fundamentalna grupa grafu grup ma prezentację
\[
BS(1,2)=\langle a,t\mid tat^{-1}=a^2\rangle.
\]
W odpowiadającym mu drzewie Bassa--Serre'a
każdy wierzchołek ma stopień trzy:
dwie orientacje pętli wnoszą odpowiednio indeksy
$[\mathbb Z:\mathbb Z]=1$ i
$[\mathbb Z:2\mathbb Z]=2$.
W przeciwieństwie do przykładu torusa
z relacją $tat^{-1}=a$,
grupa ta jest nieprzemienna. Istotnie
można zdefiniować homomorfizm do grupy
przekształceń afinicznych prostej,
przypisując $a:x\mapsto x+1$ i
$t:x\mapsto 2x$. Oznaczmy te przekształcenia przez $A,T$.
Wtedy $TAT^{-1}(x)=2(x/2+1)=x+2=A^2(x)$,
więc prezentacja wyznacza homomorfizm.
Jednak $TA(x)=2x+2$, a $AT(x)=2x+1$.
Obraz homomorfizmu jest nieprzemienny,
a zatem sama grupa także jest nieprzemienna.
\end{przyklad}

\begin{uwaga}[Literatura]
Pełną konstrukcję drzew i grafów grup opisuje J.-P. Serre,
\emph{Trees} oraz wykład J. F. Manninga
\emph{Bass--Serre Theory and Complexes of Groups}.
\end{uwaga}


% !TeX encoding = UTF-8
% Teoria wezlow -- tresc bez preambuly.
% W ksiazce: \input{teoria_wezlow_tresc.tex}
% Wersja samodzielna: kompilowac teoria_wezlow.tex.
% Wymagane srodowiska: definicja, twierdzenie, lemat, stwierdzenie,
% wniosek, przyklad, uwaga, proof; pakiety w preambule pliku glownego.

% Rysunki TikZ: xcolor, tikz oraz biblioteki arrows.meta,calc,positioning.
\definecolor{ktBlue}{HTML}{2369A8}
\definecolor{ktTeal}{HTML}{178F87}
\definecolor{ktCoral}{HTML}{D96152}
\definecolor{ktPurple}{HTML}{7755A7}
\definecolor{ktGold}{HTML}{C9932D}
\definecolor{ktInk}{HTML}{283747}


